\documentclass[11pt]{article}
\usepackage[T1]{fontenc}
\usepackage{lmodern,amsmath,amssymb,amsthm,mathtools,booktabs,microtype}
\usepackage[margin=1in]{geometry}
\usepackage{xcolor}
\usepackage[colorlinks=true,linkcolor=blue!50!black,urlcolor=blue!50!black,citecolor=blue!50!black]{hyperref}
\newtheorem{theorem}{Theorem}
\newtheorem{lemma}[theorem]{Lemma}
\newtheorem{corollary}[theorem]{Corollary}
\newtheorem{proposition}[theorem]{Proposition}
\theoremstyle{remark}\newtheorem{remark}[theorem]{Remark}
\newcommand{\R}{\mathbb R}\newcommand{\C}{\mathbb C}\newcommand{\D}{\mathbb D}
\newcommand{\arcosh}{\operatorname{arcosh}}\newcommand{\ord}{\operatorname{ord}}
\newcommand{\F}{\mathcal F}\newcommand{\dd}{\mathrm d}
\newcommand{\repo}{https://github.com/VladimirReshetnikov/ProveIt/blob/main/Analysis/GeneralizedRiemannHypothesis/}
\title{Positive filters for Dirichlet zeros:\\ a joint support--attenuation bound\\ and improved logarithmic constants}
\author{Research note for the ProveIt GRH development}
\date{27 November 2026}
\begin{document}
\maketitle
\begin{abstract}
We study positive probability measures supported on a short interval whose Fourier transforms cancel all noncentral zero ordinates of a fixed conjugation-closed family of primitive Dirichlet $L$-functions up to height $B$, with aggregate multiplicities, and are at most $e^{-L}$ at every frequency of absolute value at least $B$. A slit-plane Jensen inequality yields a joint lower bound on the half-width: the costs of the zeros and attenuation \emph{add}, rather than merely giving separate lower bounds. A cosine-window interpolation construction reduces the logarithmic upper constant for cancellation to $\kappa d$, where $d$ is the family size and $\kappa=\pi/(4y_*)\approx0.674$, with $y_*\cot y_*=1/2$. An optimal positive high-frequency filter contributes only $\arcosh(e^L)/B$. Consequently, uniformly in $L\ge0$, the half-width is between $L/B+(d/2)\log B+O_{\F}(1)$ and $L/B+(\kappa d+o_{\F}(1))\log B$. The filter theory and counting estimates are unconditional; applying them to the Fourier expansion of prime-race errors uses Dirichlet GRH, the input theorem of the repository. Neither simplicity nor independence of zero ordinates is required.
\end{abstract}

\section{Main result and relation to the source}
Let $\F$ be a fixed nonempty finite multiset of primitive Dirichlet characters, closed under conjugation with multiplicities, and let $d=|\F|$. Allow the conductor-one character, whose $L$-function is $\zeta$. Write $q_\chi$ for the conductor. For $B\ge2$ let $\Lambda_B$ be the symmetric multiset of nonzero ordinates $\gamma$ of nontrivial zeros with $|\gamma|\le B$, counted with multiplicity and added over $\F$. In this notation ordinates are real even without GRH. For a probability measure $\mu$ on $[-H,H]$ put
\[
 F(z)=\widehat\mu(z)=\int e^{izt}\,\dd\mu(t),\qquad z\in\C.
\]
Define $H_{\F}(B,L)$ as the infimum of such half-widths $H\ge0$ subject to
\begin{equation}\label{eq:conditions}
 \ord_\gamma F\ge m_{\Lambda_B}(\gamma)\quad(\gamma\in\Lambda_B),\qquad
 \sup_{|x|\ge B}|F(x)|\le e^{-L},\quad L\ge0.
\end{equation}
Zero order means complex analytic order. The same infimum is obtained with even probability densities: symmetrization preserves both conditions, and convolution with a uniform density of arbitrarily small support preserves zeros and cannot increase the real-axis supremum.

\begin{theorem}[Joint support--attenuation bounds]\label{thm:main}
There is a unique $y_*\in(0,\pi/2)$ satisfying $y_*\cot y_*=1/2$. Set $\kappa=\pi/(4y_*)$. For all sufficiently large $B$ and all $L\ge0$,
\begin{align}
 H_{\F}(B,L)&\ge \frac{L}{B}+\frac d2\log B
       +\frac12\sum_{\chi\in\F}\log\frac{q_\chi}{4\pi}-\frac d2
       -O_{\F}\!\left(\frac{\log^2 B}{B}\right),\label{eq:lower}\\
 H_{\F}(B,L)&\le \frac{\arcosh(e^L)}{B}+\kappa d\log B
       +O_{\F}\!\left(\sqrt{\log B}+(\log\log B)^2\right).\label{eq:upper}
\end{align}
The constants in the error bounds are independent of $L$. Numerically $\kappa\approx0.6738$ (see the accompanying script for additional digits).
\end{theorem}
In particular, for $L=cB\log B$ with $c\ge0$ fixed,
\begin{equation}\label{eq:leading}
 c+\frac d2\ \le\ \liminf_{B\to\infty}\frac{H_{\F}(B,cB\log B)}{\log B}
 \ \le\ \limsup_{B\to\infty}\frac{H_{\F}(B,cB\log B)}{\log B}
 \ \le\ c+\kappa d.
\end{equation}
If $L/(B\log B)\to\infty$, then $H_{\F}(B,L)\sim L/B$.

\paragraph{Source audit.} The repository's vendored \texttt{grh/src/Statements.lean} and \texttt{grh/src/Main.lean} explicitly state the Dirichlet zero-line theorem and its primitive specialization \cite{sourcelean}. The root README describes the external commit as input, while the article README describes a merge, not an independent proof audit. The article's reconstruction section proves a triangular-window interpolation lemma, a coloring construction, a logarithmic-width filter, a Jensen support bound, and a simultaneous averaging application (labels \texttt{rec:lem:sepfilter}, \texttt{rec:thm:finitefilter}, \texttt{rec:thm:logfilter}, \texttt{rec:lem:jensensupport}, \texttt{rec:thm:polylogavg}) \cite{rec}. Its final open-questions subsection asks for the optimal constant of logarithmic support. Our contribution is a joint Green-function lower bound, an improved leading cancellation constant via a growing separation scale and a cosine window, and an optimal independent stopband construction. The remaining constant gap is not closed. The classical tools themselves are not claimed as new, and a limited literature search is not proof of priority. The precise files and blob identifiers consulted are in \texttt{SOURCE\_AUDIT.md}.

\section{A slit-plane Jensen inequality}
\begin{theorem}[Weighted joint lower bound]\label{thm:slit}
Let $\mu$ be any probability measure on $[-H,H]$, and $B>0$. If $0<\lambda_j\le B$ are prescribed positive real zeros of $F$, counted with multiplicity, then
\begin{equation}\label{eq:joint}
 HB\ \ge\ 2\sum_j\arcosh\frac B{\lambda_j}
  +\frac2\pi\int_0^{\pi/2}\log\frac1{|F(B\sec\theta)|}\,\dd\theta.
\end{equation}
In particular, if $|F(x)|\le e^{-L}(B/|x|)^s$ for $|x|\ge B$, $L,s\ge0$, then
\begin{equation}\label{eq:power}
 H\ge \frac{L+s\log2}{B}+\frac2B\sum_j\arcosh\frac B{\lambda_j}.
\end{equation}
The measure need not be even. The lower bound is valid also when a prescribed zero is exactly $B$ (its displayed contribution is zero).
\end{theorem}
\begin{proof}
Compact support gives $F(0)=1$, $|F(z)|\le e^{H|\Im z|}$ and $F(-x)=\overline{F(x)}$ on the real axis. On the unit disk define
\[
 z(w)=\frac{2Bw}{1+w^2},\qquad S(w)=\frac{B(1-w^2)}{1+w^2},\qquad
 G(w)=F(z(w))e^{-HS(w)}.
\]
These are analytic on $\D$, with $S(w)^2=B^2-z(w)^2$. For $w=x+iy$, $r=|w|<1$,
\[
 \Re S(w)=\frac{B(1-r^4)}{|1+w^2|^2},\qquad
 |\Im z(w)|=\frac{2B|y|(1-r^2)}{|1+w^2|^2}.
\]
As $1+r^2\ge2|y|$, $|G|\le1$ on $\D$ and $G(0)=e^{-HB}$. For $0<\lambda<B$ the solution of $z(w)=\lambda$ in $\D$ is
\[
 w_\lambda=\frac{\lambda}{B+\sqrt{B^2-\lambda^2}},\qquad
 -\log w_\lambda=\arcosh(B/\lambda).
\]
Its negative corresponds to $-\lambda$; the map has nonzero derivative at each of these points, so analytic multiplicities are preserved. Conjugation of $F$ shows both signs have equal orders. Jensen's formula at any radius $r<1$ not passing through a zero gives
\[
 HB=\sum_{|a|<r:\,G(a)=0}\log\frac r{|a|}
       +\frac1{2\pi}\int_0^{2\pi}-\log|G(re^{i\theta})|\,\dd\theta,
\]
where zeros are repeated. Retain the prescribed interior zeros and take $r\uparrow1$. The integrand is nonnegative, so Fatou's lemma applies, even at boundary zeros. Except at $w=\pm i$, the boundary values satisfy $z(e^{i\theta})=B\sec\theta$ and $\Re S(e^{i\theta})=0$. Evenness of the modulus and reflection of the angular variable now give \eqref{eq:joint}. The integral is automatically finite by this inequality. Finally, $\log(1/|F(B\sec\theta)|)\ge L+s\log\sec\theta$ and $\int_0^{\pi/2}\log\sec\theta\,\dd\theta=(\pi/2)\log2$, proving \eqref{eq:power}.
\end{proof}

\paragraph{Why the two terms matter.} The circular Jensen estimate in \cite{rec} gives $HB\ge\pi\sum_j\log(B/\lambda_j)$, independently of $L$. The separate no-zero bound $HB\ge L$ does not imply their sum. Equation \eqref{eq:joint} accounts simultaneously for attenuation on the two rays and the prescribed interior zero divisor. For an isolated very low zero the older circular bound can be stronger; neither should be discarded. The weight in \eqref{eq:joint} is harmonic measure on the rays, not a uniform frequency average.

\section{The arithmetic cost of the zeros}
Write $N(t)=\sum_{0<\gamma\le t}m_{\Lambda_\infty}(\gamma)$. The classical zero count (for conductor $>1$, see \cite{count}; for $\zeta$, Riemann--von Mangoldt \cite{mv}) and conjugation closure give
\begin{equation}\label{eq:count}
 N(t)=\frac{t}{2\pi}\sum_{\chi\in\F}\left(\log\frac{q_\chi t}{2\pi}-1\right)+O_{\F}(\log(t+2)),\qquad t\ge2.
\end{equation}
Central zeros are excluded and only change the bounded part of the count. There is a least positive ordinate in the family; hence $N(t)=0$ near zero. Stieltjes integration by parts gives
\begin{equation}\label{eq:stieltjes}
 \sum_{0<\gamma\le B}m(\gamma)\arcosh(B/\gamma)
   =\int_0^B\frac{B N(t)}{t\sqrt{B^2-t^2}}\,\dd t.
\end{equation}
This also follows by integrating the derivative of $\arcosh(B/t)$ from each ordinate to $B$ and summing. Inserting \eqref{eq:count} on $[2,B]$ and extending its main term down to zero costs $O_{\F}(1)$: the main term is $O_{\F}(t(1+|\log t|))$ there, and the true count vanishes near zero. The error on $[2,B/2]$ is $O_{\F}(\log^2 B)$; on $[B/2,B]$ it is $O_{\F}(\log B)$ by the integrable square-root singularity. With $t=B\sin\theta$ and the classical integral of $\log\sin\theta$ we obtain
\begin{equation}\label{eq:sum}
 \frac2B\sum_{0<\gamma\le B}m(\gamma)\arcosh(B/\gamma)
 =\frac d2\log B+\frac12\sum_{\chi\in\F}\log\frac{q_\chi}{4\pi}-\frac d2
  +O_{\F}(\log^2 B/B).
\end{equation}
Thus \eqref{eq:power} proves \eqref{eq:lower}. For comparison the article's circular Jensen bound, evaluated from the same count, is $\tfrac d2\log B+\tfrac12\sum_\chi\log(q_\chi/(2\pi))-d+O_{\F}(\log^2 B/B)$. The slit bound improves that zero-only constant by $d(1-\log2)/2$, and, more importantly, adds $L/B$.

\section{Improved positive interpolation}
The next lemma sharpens the article's triangular-window separated interpolation. It uses a nonnegative cosine taper and does not impose any restrictions on phases or multiplicities beyond separation within a single color.
\begin{lemma}[Cosine-window interpolation]\label{lem:cosine}
Suppose $\Gamma\subset\R\setminus\{0\}$ is finite and symmetric, and all distinct points of $\Gamma\cup\{0\}$ are at distance at least $\delta>0$. There is an even probability density supported on $[-A,A]$, where
\[
 A=\frac{\pi^2}{2y_*\delta},\qquad y_*\cot y_*=\frac12,
\]
whose Fourier transform vanishes on $\Gamma$.
\end{lemma}
\begin{proof}
The function $y\cot y$ decreases strictly from $1$ to $0$ on $(0,\pi/2)$, since its derivative is $(\sin y\cos y-y)/\sin^2y<0$. Let $a=\pi/(2A)$, $u=a/\delta=y_*/\pi<1/2$. The even density
\[
 g_A(t)=\frac a2\cos(at)\,\mathbf 1_{[-A,A]}(t)
\]
has transform $f_A(x)=a^2\cos(Ax)/(a^2-x^2)$, with removable values at $x=\pm a$. In particular $f_A(0)=1$, and $|f_A(x)|\le a^2/(x^2-a^2)$ for $|x|\ge\delta$. Put
\[
 M_{\lambda\eta}=f_A(\lambda-\eta),\qquad b_\lambda=f_A(\lambda),\qquad \lambda,\eta\in\Gamma.
\]
On either side of any point of $\Gamma\cup\{0\}$, the $j$-th closest point is at distance at least $j\delta$. Therefore the induced column-sum norm and vector $\ell^1$ norm satisfy
\[
 \|M-I\|_1,\ \|b\|_1\le2\sum_{j\ge1}\frac{u^2}{j^2-u^2}
   =1-\pi u\cot(\pi u)=\frac12.
\]
The series identity follows by logarithmic differentiation of Euler's sine product $\sin(\pi u)/(\pi u)=\prod_{j\ge1}(1-u^2/j^2)$; convergence is absolute at the chosen $u$. Hence $M$ is invertible, and $c=-M^{-1}b$ satisfies $\|c\|_1\le1$. All coefficients are real and $c_{-\lambda}=c_\lambda$ by reflection symmetry and uniqueness. Set
\[
 Q(t)=1+\sum_{\lambda\in\Gamma}c_\lambda e^{-i\lambda t}.
\]
Then $Q$ is real and even, and $Q\ge0$ because $\|c\|_1\le1$. Although equality is allowed, $Q$ cannot vanish identically on an interval: distinct real exponentials are linearly independent, as follows from the Vandermonde matrix of derivatives at any interior point. Thus $Z=\int g_AQ>0$. The normalized density $g_AQ/Z$ is nonnegative, even, and its transform at $\eta\in\Gamma$ is $(b_\eta+(Mc)_\eta)/Z=0$.
\end{proof}

\begin{proposition}[Growing separation scale]\label{prop:growing}
For the family $\F$, there is an even positive probability measure $\nu_B$ cancelling $\Lambda_B$ with its aggregate multiplicities, of half-width at most
\[
 \kappa d\log B+O_{\F}(\sqrt{\log B}+(\log\log B)^2).
\]
For large $B$ it can be chosen to be a density.
\end{proposition}
\begin{proof}
Take $\delta=\sqrt{\log B}$ for large $B$. For each positive occurrence $0<\gamma\le\delta$, convolve a symmetric two-atom measure on $\{\pm\pi/(2\gamma)\}$; its transform $\cos(\pi z/(2\gamma))$ has a simple zero at both signs of $\gamma$. The sum of these half-widths is $O_{\F}(1+\log^2(\delta+2))$ by partial summation of \eqref{eq:count} and the least-positive-ordinate bound.

For the remaining positive occurrences in $(\delta,B]$, let $K$ be the maximum number (with repetitions) lying in an interval $(x,x+\delta]$, $x\ge0$. Color in nondecreasing order using $K$ colors, reusing a color only if its last point is at distance at least $\delta$ from the new point. At most $K-1$ previous occurrences can block colors, by the definition of $K$ and because the current point is $>\delta$. Thus every color class, reflected in zero, satisfies the separation lemma including separation from zero. In each color use Lemma~\ref{lem:cosine}, then convolve all the measures. Orders of vanishing add, even if the same ordinate occurs in different characters or with higher multiplicity.

To estimate $K$, extend the main term of \eqref{eq:count} to $t=0$ by continuity and enlarge the constant to make the error bound valid on $[0,2]$. For a window meeting $(\delta,B]$, truncate at $B$ if necessary. The main-term derivative is $(2\pi)^{-1}\sum_\chi\log(q_\chi t/(2\pi))\le (d/(2\pi))\log B+O_{\F}(1)$ for $0<t\le B$. The two count errors together are $O_{\F}(\log B)$. Hence
\[
 K\le \frac{d\delta}{2\pi}\log B+O_{\F}(\delta+\log B).
\]
The high-part half-width is at most $K\pi^2/(2y_*\delta)$, which equals at most $\kappa d\log B+O_{\F}(1+\log B/\delta)$. Add the low-part cost and substitute the chosen $\delta$. At least one high color is nonempty for all sufficiently large $B$ by \eqref{eq:count}, so the convolution is a density.
\end{proof}

\section{Optimal uniform suppression without prescribed zeros}
The following positive kernel is a continuous Chebyshev-type window; we prove its exact minimax property directly. Classical Chebyshev window design is prior art \cite{cheb}. The separate minimax result is used as a component, not asserted as a new classical extremal theorem.
\begin{theorem}[Exact stopband minimax]\label{thm:stop}
For $B>0$ and $h\ge0$,
\[
 \inf_{\substack{\mu\text{ probability}\operatorname{supp}\mu\subset[-h,h]}}
  \sup_{|x|\ge B}|\widehat\mu(x)|=\frac1{\cosh(Bh)}.
\]
An extremizer has entire transform
\begin{equation}\label{eq:cheb}
 C_{B,h}(z)=\frac{\cosh(h\sqrt{B^2-z^2})}{\cosh(Bh)}.
\end{equation}
The square root notation denotes the entire power series in $B^2-z^2$.
\end{theorem}
\begin{proof}
For positivity of \eqref{eq:cheb}, consider a continuous-time particle starting with velocity $+1$ or $-1$ equally likely, reversing velocity at each event of a Poisson process of rate $B$. Condition on an even number of reversals up to time $h$. Its displacement $X$ lies in $[-h,h]$ and has a symmetric distribution. To compute its transform, the matrix of transforms with specified starting and ending velocities is
\[
 e^{-Bh}\exp\left(h\begin{pmatrix} iz&B\\ B&-iz\end{pmatrix}\right).
\]
Indeed the small-time transition has diagonal term $1-B\,\dd t\pm iz\,\dd t$ and off-diagonal term $B\,\dd t$, giving the constant-coefficient differential equation; equivalently one expands over ordered reversal times. Even reversals mean equal starting and ending velocities, so the unnormalized transform is half the trace of this matrix. Its square inside the exponential is $(B^2-z^2)I$, whence half the trace is $e^{-Bh}\cosh(h\sqrt{B^2-z^2})$. At zero this is the probability of even reversals, $e^{-Bh}\cosh(Bh)>0$. Dividing gives \eqref{eq:cheb}. For $|x|\ge B$ the numerator is $\cos(h\sqrt{x^2-B^2})$; its largest absolute value is $1$ (already at $x=B$).

For the lower bound let $\epsilon=\sup_{|x|\ge B}|F(x)|$. It cannot be zero, by the identity theorem. Symmetrize the measure; its transform is then real on the real axis and still bounded there by $\epsilon$ on the stopband. Let $v(\xi)=\log|\xi+\sqrt{\xi^2-1}|$ off $[-1,1]$, choosing the root so that the modulus is $>1$, and set $v=0$ on $[-1,1]$. This is a continuous subharmonic function, with $v(\xi)\le\log(2|\xi|+1)$. For example, the logarithmic-potential representation
\[
 v(\xi)=\frac1\pi\int_0^\pi\log|2(\xi-\cos\theta)|\,\dd\theta
\]
follows from the Joukowski map and the elementary circle mean for $\log|z-a|$, and gives subharmonicity (or use the harmonic Green function extended by zero). Use the disk maps $z(w),S(w)$ from Theorem~\ref{thm:slit} and set
\[
 U(w)=v(F(z(w))/\epsilon)-h\Re S(w).
\]
This is subharmonic, bounded above by $\log(2/\epsilon+1)$ because $\Re S\ge|\Im z|$. At every boundary point except $\pm i$, its radial limit is zero: $F(z(w))/\epsilon$ is real in $[-1,1]$ there. The subharmonic mean inequality at the origin, followed by reverse Fatou (using the common upper bound), gives $U(0)\le0$. Since $F(0)=1$, this says $\arcosh(1/\epsilon)\le Bh$. Thus $\epsilon\ge\operatorname{sech}(Bh)$ for the original measure as well. The case $h=0$ is immediate.
\end{proof}

\paragraph{Completion of Theorem~\ref{thm:main}.} Convolve $\nu_B$ of Proposition~\ref{prop:growing} with the measure in Theorem~\ref{thm:stop} at $h=B^{-1}\arcosh(e^L)$. Supports add, the zero orders persist, and the second factor makes the tail at most $e^{-L}$ because the first is a characteristic function of modulus at most one. This is \eqref{eq:upper}. Since $L\le\arcosh(e^L)\le L+\log2$, \eqref{eq:leading} and the strong-attenuation asymptotic follow.

\section{Consequences for the repository's GRH program}
\subsection*{A sharper simultaneous prime-race averaging rule (uses GRH)}
Fix $q\ge2$, let $n_q=\varphi(q)$ and take the family of primitive characters inducing the characters modulo $q$ (one entry per character). It is conjugation-closed and $d=n_q$. In the notation of the article's ramp formula (\texttt{zer:lem:ramp}, \texttt{zer:eq:Zexp}) \cite{zer}, define
\begin{align*}
 S_a(x)&=\sum_{n\le x,\ n\equiv a\ (q)}\Lambda(n)(1/n-1/x),\\
 R_a(x)&=S_a(x)-\log x/n_q-c_a,\\
 e^{t/2}R_a(e^t)&=\beta_a+\sum_{\gamma\ne0}c_{a,\gamma}e^{i\gamma t}
       +O_q((1+t)e^{-t/2}),\qquad t\ge\log2.
\end{align*}
Here $\Lambda$ is the von Mangoldt function, $c_a$ is the constant term of the ramp expansion, $\beta_a=4n_q^{-1}\sum_{\chi\bmod q}\overline{\chi(a)}m_\chi(0)$, and $c_{a,\gamma}=n_q^{-1}\sum_{\chi\bmod q}\overline{\chi(a)}m_\chi(\gamma)/(1/4+\gamma^2)$. GRH makes the zero terms pure Fourier frequencies. The local zero count gives absolute convergence and the tail estimate
\[
 \sum_{|\gamma|>B}|c_{a,\gamma}|=O_q(\log(B+2)/B),
\]
by partial summation. No independence of the ordinates is involved. For clarity, the ramp expansion follows by Mellin inversion of $y^s/(s(s-1))$ against $-L'/L$, moving the contour from $\Re s>1$ to $-1/2$. The pole at $1$ gives the main and constant terms, the residue at a zero $\rho=1/2+i\gamma$ is $m_\rho x^{\rho-1}/(\rho(1-\rho))$, and the possible pole at $0$ contributes $O_q((1+\log x)/x)$. Removed Euler factors of imprimitive characters contribute constants plus the same error; character orthogonality yields the residue-class expansion. This is the analytic input from the article, not a new derivation of GRH.

\begin{corollary}[Explicit leading polylogarithmic range]\label{cor:prime}
Assume Dirichlet GRH. Fix any $a_0>0$ and $\eta>0$. For sufficiently large $X$, there is a common even probability density $\mu_X$ supported on $[-H,H]$ with
\[
 H\le(\kappa\varphi(q)+\eta)\log\log X,
\]
such that simultaneously for all reduced classes $a$,
\begin{equation}\label{eq:prime}
 \int e^{u/2}R_a(Xe^u)\,\dd\mu_X(u)=\frac{\beta_a}{\sqrt X}
 +O_{q,a_0,\eta}\left(\frac{(\log X)^{1+(\kappa\varphi(q)+\eta)/2}}{X}+X^{-a_0-1/2}\right).
\end{equation}
The kernel is independent of the residue class, has no negative weights, and uses only scales between $X(\log X)^{-\kappa\varphi(q)-\eta}$ and $X(\log X)^{\kappa\varphi(q)+\eta}$.
\end{corollary}
\begin{proof}
Set $T=\log X$, $B=T$, $L=a_0T$. Theorem~\ref{thm:main} supplies a filter with $H\le(\kappa n_q+\eta)\log T$ for all sufficiently large $T$, as $L/B=a_0$ is constant. Average the absolutely convergent Fourier series at $t=T+u$. Its low frequencies vanish, its high frequencies have total contribution $O_q(e^{-L}\log B/B)=O_q(e^{-L})$, and the analytic remainder is bounded by $O_q((1+T+H)e^{-(T-H)/2})$ since $T-H\to\infty$. Divide the result by $\sqrt X$. This gives \eqref{eq:prime}. In particular $a_0\ge1/2$ makes the high-frequency error no larger than $X^{-1}$.
\end{proof}

For the exact-cancellation and full-stopband requirements used here, the lower bound forces $H\ge(d/2)\log\log X+O_q(1)$, even if only constant $a_0T$ attenuation is requested. This is not a lower bound for every possible prime-race averaging algorithm: cancellation only of distinct nonzero Fourier modes, center-dependent quadrature, or direct cancellation of the weighted error can be easier. The aggregate multiplicities are an explicit design requirement in this note, not an asserted necessity for prime-race applications.

\subsection*{A numerical comparison of leading constants}
For one zeta spectrum ($d=1$), let $c_0=-\log(\sin1)$, the tail exponent used in the article's sinc-power construction. The rigorous asymptotic coefficient brackets are as follows (rounded for illustration).
\begin{center}
\begin{tabular}{@{}lrr@{}}
\toprule
Tail requirement $L$ & lower coefficient & upper coefficient\\\midrule
$0$ & $0.5000$ & $0.6738$\\
$c_0 B\log B$ & $0.6726$ & $0.8464$\\
$B\log B$ & $1.5000$ & $1.6738$\\
\bottomrule
\end{tabular}
\end{center}
The error $o(\log B)$ in these coefficient statements is for fixed $\F$. For finite $B$ one should instead use \eqref{eq:joint} and the explicit coloring construction. The provided code computes the constants, checks the transform formula, numerically evaluates example finite zero lists, and verifies the interpolation residual for a clustered example. It is not a zero certification or a formal proof checker.

\section{Further research questions}
\begin{enumerate}
\item Is the exact cancellation constant $d/2$ attainable, or does positivity force a strictly larger constant? Can one improve the cosine interpolation or use an entirely different positive kernel? The lower and upper leading constants above leave a gap of $(\kappa-1/2)d$.
\item For a finite prescribed divisor, characterize the optimizer for the joint problem, including possible atoms at the support endpoints. Can the weighted Jensen inequality be strengthened to an exact constrained minimax formula analogous to Theorem~\ref{thm:stop}?
\item Can one replace aggregate multiplicities by distinct ordinates and derive sharp constants without assumptions on common zeros? This would more closely match the requirements of ordinary prime-race Fourier cancellation.
\item Can the positive interpolation construction be made computationally certified when the target ordinates are given by intervals rather than exact values? The main difficulty is controlling residuals and width while handling clusters and multiplicities.
\end{enumerate}

\paragraph{Verification status.} All new assertions above have mathematical proofs, using standard analytic facts (Jensen's formula, subharmonicity, the classical zero count, and the ramp explicit formula). The source's Lean theorem statements were inspected; the vendored project was not rebuilt or independently verified here, and this note is not Lean-formalized. Statements about arithmetic prime errors explicitly depend on Dirichlet GRH; the new filter inequalities do not. The repository was not modified. No claim of resolving a further named conjecture is made.
\begin{thebibliography}{9}
\bibitem{sourcelean} ProveIt contributors, \emph{Vendored grh-lean}, upstream commit \texttt{72fead102f663d5c7a407964a13dc0cf9d977207}. \href{\repo vendor/grh-lean/grh/src/Main.lean}{Main.lean}; \href{\repo vendor/grh-lean/grh/src/Statements.lean}{Statements.lean}.
\bibitem{rec} ProveIt contributors, \emph{Beyond the Riemann hypothesis for Dirichlet L-functions: consequences, transfers and limits}, \href{\repo Article/sections/rec.tex}{Recovering zeros and characters from prime data}, section on shortest positive averages; accessed 27 November 2026.
\bibitem{zer} Same article, \href{\repo Article/sections/zer.tex}{Zeros, Euler products and prime races under Dirichlet GRH}, ramp explicit formula.
\bibitem{count} M. A. Bennett, G. Martin, K. O'Bryant and A. Rechnitzer, Counting zeros of Dirichlet $L$-functions, \emph{Math. Comp.} 90 (2021), 1455--1482. \href{https://doi.org/10.1090/mcom/3599}{doi:10.1090/mcom/3599}.
\bibitem{mv} H. L. Montgomery and R. C. Vaughan, \emph{Multiplicative Number Theory I: Classical Theory}, Cambridge University Press, 2007, Ch.~14 (Riemann--von Mangoldt formula).
\bibitem{cheb} P. Lynch, The Dolph--Chebyshev window: a simple optimal filter, \emph{Monthly Weather Review} 125 (1997), 655--660. \href{https://journals.ametsoc.org/doi/10.1175/1520-0493(1997)125\%3C0655\%3ATDCWAS\%3E2.0.CO\%3B2}{Journal reference}. This is background on classical Chebyshev windows, not a source for the joint inequality.
\end{thebibliography}
\end{document}
